\chapter{Complexity — P Subset P Poly Config}
%%% generated by the blueprint pipeline from TCSlib.Complexity.CircuitComplexity.PSubsetPPolyConfig
\section{Overview}

Declarations of the Lean module \texttt{TCSlib.Complexity.CircuitComplexity.PSubsetPPolyConfig}.

\section{Declarations}

\begin{definition}[Configuration width]
\lean{Complexity.cfgWidth}
\label{Complexity.cfgWidth}
\leanok
\uses{Complexity.snapWidth}
For a Boolean-input finite Turing machine $M$, define its configuration width to be the
snapshot width of $M$ plus $5$.
\end{definition}

\begin{definition}[Configuration arity]
\lean{Complexity.cfgArity}
\label{Complexity.cfgArity}
\leanok
\uses{Complexity.snapWidth}
For a Boolean-input finite Turing machine $M$, define its configuration arity to be the
snapshot width of $M$ plus twice its parameter $k$, plus $13$.
\end{definition}

\begin{definition}[Configuration stride]
\lean{Complexity.cfgStride}
\label{Complexity.cfgStride}
\leanok
For a Boolean-input finite Turing machine $M$ and natural numbers $N$ and $T$, define
the configuration stride to be $k(2T+1)+N+3$, where $k$ is the parameter of $M$.
\end{definition}

\begin{definition}[Configuration cell index]
\lean{Complexity.cellIdx}
\label{Complexity.cellIdx}
\leanok
\uses{Complexity.cfgStride}
For a Boolean-input finite Turing machine $M$, natural numbers $N,T,t,r$, and an index
$\tau$ with $0\leq\tau<k$, define the cell index by $t$ times the configuration stride,
plus $\tau(2T+1)+r$. Here the configuration stride is $k(2T+1)+N+3$.
\end{definition}

\begin{definition}[Configuration input index]
\lean{Complexity.inpIdx}
\label{Complexity.inpIdx}
\leanok
\uses{Complexity.cfgStride}
For a Boolean-input finite Turing machine $M$, natural numbers $N,T,t,p$, define the
input index to be $t$ times the configuration stride, plus $k(2T+1)+p$, where the
configuration stride is $k(2T+1)+N+3$.
\end{definition}

\begin{definition}[Configuration snapshot index]
\lean{Complexity.snapIdx}
\label{Complexity.snapIdx}
\leanok
\uses{Complexity.cfgStride}
For a Boolean-input finite Turing machine $M$, natural numbers $N,T,t$, define the
snapshot index to be $t$ times the configuration stride, plus $k(2T+1)+N+2$, where the
configuration stride is $k(2T+1)+N+3$.
\end{definition}

\begin{definition}[Configuration positions]
\lean{Complexity.CfgPos}
\label{Complexity.CfgPos}
\leanok
For a natural number $k$, configuration positions are of three kinds: a work-cell
position, specified by a time $t$, a track $\tau \text{ with } 0 \leq \tau < k$, and a
cell coordinate $r$; an input position, specified by a time $t$ and an input coordinate
$p$; or a snapshot position, specified by a time $t$.
\end{definition}

\begin{definition}[Indexing configuration positions]
\lean{Complexity.cfgPos}
\label{Complexity.cfgPos}
\leanok
\uses{Complexity.CfgPos, Complexity.cfgStride}
Let $M$ be a finite Turing machine with Boolean tape alphabet, let $N,T$ be natural
numbers, and let $i$ be a natural number. Write $s$ for the stride determined by
$M,N,T$, and write $q=\lfloor i/s\rfloor$ and $u=i \bmod s$. The position assigned to
$i$ is the work-cell position $(q,\lfloor u/(2T+1)\rfloor,u\bmod(2T+1))$ if $u<k(2T+1)$,
where $k$ is the number of work tracks of $M$; otherwise it is the input position
$(q,u-k(2T+1))$ if $u<k(2T+1)+N+2$, and the snapshot position at time $q$ in all other
cases.
\end{definition}

\begin{theorem}[Positivity of the configuration stride]
\lean{Complexity.stride_pos}
\label{Complexity.stride_pos}
\leanok
\uses{Complexity.cfgStride}
\difficulty{2}
Let $M$ be a finite-state Turing machine with $k$ work tapes and Boolean tape symbols,
and let $N,T\in\mathbb{N}$. The stride $k(2T+1)+N+3$ is positive.
\end{theorem}

\begin{theorem}[Quotient and remainder of a stride encoding]
\lean{Complexity.div_mod_stride}
\label{Complexity.div_mod_stride}
\leanok
\uses{Complexity.cfgStride, Complexity.stride_pos}
\difficulty{3}
Let $M$ be a Turing machine with $k(M)$ states, and let $N,T,t,o$ be natural numbers.
Define the stride $s=k(M)(2T+1)+N+3$. If $o<s$, then dividing $ts+o$ by $s$ gives
quotient $t$ and remainder $o$.
\end{theorem}

\begin{theorem}[Cell offset bound]
\lean{Complexity.cell_off_lt}
\label{Complexity.cell_off_lt}
\leanok
\difficulty{3}
Let $M$ be a finite Turing machine with $k$ work tapes, let $N,T\in\mathbb{N}$, and let
$\tau<k$ and $r<2T+1$. Then $\tau(2T+1)+r<k(2T+1)$.
\end{theorem}

\begin{theorem}[Decoding work-cell indices]
\lean{Complexity.cfgPos_cellIdx}
\label{Complexity.cfgPos_cellIdx}
\leanok
\uses{Complexity.cellIdx, Complexity.cell_off_lt, Complexity.cfgPos, Complexity.cfgStride, Complexity.div_mod_stride}
\difficulty{5}
Let $M$ be a finite-state binary Turing machine with $k$ work tapes, let
$N,T,t\in\mathbb N$, and let $\tau$ be a work tape. For every offset $r<2T+1$, the
tableau index assigned to time $t$, tape $\tau$, and offset $r$ decodes as the work-cell
position $(t,\tau,r)$.
\end{theorem}

\begin{theorem}[Input position indexing]
\lean{Complexity.cfgPos_inpIdx}
\label{Complexity.cfgPos_inpIdx}
\leanok
\uses{Complexity.cfgPos, Complexity.cfgStride, Complexity.div_mod_stride, Complexity.inpIdx}
\difficulty{4}
Let $M$ be a finite-state Turing machine with $k$ work tapes, and let $N,T,t,p$ be
natural numbers with $p<N+2$. Set $S=k(2T+1)+N+3$. The configuration position encoded by
the index $tS+k(2T+1)+p$ is the input position $(t,p)$.
\end{theorem}

\begin{theorem}[Snapshot indexing]
\lean{Complexity.cfgPos_snapIdx}
\label{Complexity.cfgPos_snapIdx}
\leanok
\uses{Complexity.cfgPos, Complexity.cfgStride, Complexity.div_mod_stride, Complexity.snapIdx}
\difficulty{4}
Let $M$ be a finite-state Turing machine with $k$ work tapes, and let $N,T,t$ be natural
numbers. Set $S=k(2T+1)+N+3$. The configuration position encoded by $tS+k(2T+1)+N+2$ is
the snapshot at time $t$.
\end{theorem}

\begin{theorem}[Exhaustive tableau index cases]
\lean{Complexity.cfgIdx_cases}
\label{Complexity.cfgIdx_cases}
\leanok
\uses{Complexity.cellIdx, Complexity.cfgStride, Complexity.inpIdx, Complexity.snapIdx, Complexity.stride_pos}
\difficulty{5}
Let $M$ be a finite-state binary Turing machine with $k$ work tapes, and let
$N,T\in\mathbb N$. Every natural number $i<(T+1)\bigl(k(2T+1)+N+3\bigr)$ has at least
one of the following forms: it indexes a work-tape cell at some time $t\le T$, for a
tape $\tau$ and offset $r<2T+1$; it indexes an input position at some time $t\le T$,
with position $p<N+2$; or it indexes a snapshot at some time $t\le T$.
\end{theorem}

\begin{theorem}[Cell index bound]
\lean{Complexity.cellIdx_lt_succ}
\label{Complexity.cellIdx_lt_succ}
\leanok
\uses{Complexity.cellIdx, Complexity.cell_off_lt, Complexity.cfgStride}
\difficulty{3}
Let $M$ be a finite Turing machine with $k$ work tapes, and let $N,T,t\in\mathbb{N}$,
$\tau<k$, and $r<2T+1$. Define the configuration stride as $k(2T+1)+N+3$ and the cell
index as $t$ times this stride plus $\tau(2T+1)+r$. Then the cell index is less than
$(t+1)$ times the stride.
\end{theorem}

\begin{theorem}[Input index precedes the next block]
\lean{Complexity.inpIdx_lt_succ}
\label{Complexity.inpIdx_lt_succ}
\leanok
\uses{Complexity.cfgStride, Complexity.inpIdx}
\difficulty{3}
Let $M$ be a finite-state multi-tape Turing machine, and let $N,T,t\in\mathbb N$. For
every $p\in\mathbb N$ with $p<N+2$, the index $t\bigl(k(2T+1)+N+3\bigr)+k(2T+1)+p$ is
strictly less than $(t+1)\bigl(k(2T+1)+N+3\bigr)$, where $k$ is the number of work tapes
of $M$.
\end{theorem}

\begin{theorem}[Snapshot index bound]
\lean{Complexity.snapIdx_lt_succ}
\label{Complexity.snapIdx_lt_succ}
\leanok
\uses{Complexity.cfgStride, Complexity.snapIdx}
\difficulty{3}
Let $M$ be a finite Turing machine over the Boolean alphabet, and let $N,T,t\in\bbn$.
Write $k$ for the number of work tapes of $M$. Then
$t\bigl(k(2T+1)+N+3\bigr)+k(2T+1)+N+2 < (t+1)\bigl(k(2T+1)+N+3\bigr)$.
\end{theorem}

\begin{theorem}[Lower bound for a cell index]
\lean{Complexity.le_cellIdx}
\label{Complexity.le_cellIdx}
\leanok
\uses{Complexity.cellIdx, Complexity.cfgStride}
\difficulty{3}
Let $M$ be a finite Turing machine with $k$ work tapes, and let $N,T,t,r$ be nonnegative
integers and $\tau$ a work-tape index. If $S=k(2T+1)+N+3$ is the configuration stride
and $I$ is the cell index at time $t$, tape $\tau$, and position $r$, then $tS\le I$.
\end{theorem}

\begin{theorem}[Input index precedes snapshot index]
\lean{Complexity.inpIdx_lt_snapIdx}
\label{Complexity.inpIdx_lt_snapIdx}
\leanok
\uses{Complexity.inpIdx, Complexity.snapIdx}
\difficulty{3}
Let $M$ be a finite-state multi-tape Turing machine, and let $N,T,t\in\mathbb N$. For
every $p\in\mathbb N$ with $p<N+2$, the index $t\bigl(k(2T+1)+N+3\bigr)+k(2T+1)+p$ is
strictly less than $t\bigl(k(2T+1)+N+3\bigr)+k(2T+1)+N+2$, where $k$ is the number of
work tapes of $M$.
\end{theorem}

\begin{theorem}[Cell indices precede input indices]
\lean{Complexity.cellIdx_lt_inpIdx}
\label{Complexity.cellIdx_lt_inpIdx}
\leanok
\uses{Complexity.cellIdx, Complexity.cell_off_lt, Complexity.inpIdx}
\difficulty{3}
Let $M$ be a finite-state Turing machine over the Boolean alphabet, let $N,T$ be natural
numbers, and write $k$ for the number of states of $M$. Define the stride
$s=k(2T+1)+N+3$. For every time $t$, state index $\tau$ with $0\leq\tau<k$, row index
$r<2T+1$, and natural number $p$, the index $ts+\tau(2T+1)+r$ is strictly less than
$ts+k(2T+1)+p$.
\end{theorem}

\begin{definition}[Sources for the previous snapshot]
\lean{Complexity.prevSnapSrcs}
\label{Complexity.prevSnapSrcs}
\leanok
\uses{BoolCircuit.BitSrc, Complexity.snapIdx, Complexity.snapWidth}
Let $M$ be a finite Turing machine with Boolean tape alphabet, let $\ell$ be a finite
list of bit sources, and let $T,t$ be natural numbers. The sources for the snapshot at
time $t$ form a list of length equal to the snapshot width: if $t=0$, every entry is the
constant false bit; otherwise, entry $j$ is bit $j$ of the block computed by the
instruction for the snapshot at time $t-1$.
\end{definition}

\begin{definition}[Neighbour source for a work cell]
\lean{Complexity.nbSrc}
\label{Complexity.nbSrc}
\leanok
\uses{BoolCircuit.BitSrc, Complexity.cellIdx}
Let $M$ be a finite Turing machine with Boolean tape alphabet, let $\ell$ be a finite
list of bit sources, and let $T,t,r,j$ be natural numbers, with $\tau$ a work-track
index and $ok$ a Boolean. The neighbour source is bit $j$ of the block for the cell at
time $t-1$, track $\tau$, and coordinate $r$, if $t\ne0$ and $ok$ is true; otherwise it
is the constant false bit.
\end{definition}

\begin{definition}[Left-neighbour source for a work cell]
\lean{Complexity.chSrc}
\label{Complexity.chSrc}
\leanok
\uses{BoolCircuit.BitSrc, Complexity.cellIdx}
Let $M$ be a finite Turing machine with Boolean tape alphabet, let $\ell$ be a finite
list of bit sources, and let $T,t,r,j$ be natural numbers, with $\tau$ a work-track
index. The source is bit $j$ of the block for the cell at time $t$, track $\tau$, and
coordinate $r-1$, if $r\ne0$; otherwise it is the constant false bit.
\end{definition}

\begin{definition}[Sources for a tableau cell]
\lean{Complexity.cellSrcs}
\label{Complexity.cellSrcs}
\leanok
\uses{BoolCircuit.BitSrc, Complexity.chSrc, Complexity.nbSrc, Complexity.prevSnapSrcs}
Let $M$ be a Turing machine with Boolean tape alphabet, let $\ell$ be a list of bit
sources, let $T,t,r\in\mathbb{N}$, and let $\tau$ be a track index of $M$. The
cell-source list at $(t,\tau,r)$ is formed, in order, from the constants indicating
whether $t=0$ and whether $r=T$; the three neighbor sources at coordinate $r-1$ with
boundary flag $1\le r$; the three at coordinate $r$ with flag true; the three at
coordinate $r+1$ with flag $r+1<2T+1$; two cell sources at coordinate $r$ indexed by $3$
and $4$; and the previous-snapshot sources for time $t$.
\end{definition}

\begin{definition}[Source for an input-position block]
\lean{Complexity.ihSrc}
\label{Complexity.ihSrc}
\leanok
\uses{BoolCircuit.BitSrc, Complexity.inpIdx}
Let $M$ be a finite Turing machine with Boolean tape alphabet, let $\ell$ be a finite
list of bit sources, and let $T,t,p'$ be natural numbers, with $ok$ a Boolean. The
source is bit $0$ of the block for the input position $p'$ at time $t-1$ if $ok$ is
true; otherwise it is the constant false bit.
\end{definition}

\begin{definition}[Source for a preceding input position]
\lean{Complexity.ichSrc}
\label{Complexity.ichSrc}
\leanok
\uses{BoolCircuit.BitSrc, Complexity.inpIdx}
Let $M$ be a finite Turing machine with Boolean tape alphabet, let $\ell$ be a finite
list of bit sources, and let $T,t,p,j$ be natural numbers. The source is bit $j$ of the
block for input position $p-1$ at time $t$ if $p\ne0$; otherwise it is the constant
false bit.
\end{definition}

\begin{definition}[Input source list]
\lean{Complexity.inpSrcs}
\label{Complexity.inpSrcs}
\leanok
\uses{BoolCircuit.BitSrc, Complexity.ichSrc, Complexity.ihSrc, Complexity.prevSnapSrcs}
Let $M$ be a Turing machine with Boolean symbols, let $\ell$ be a finite list of bit
sources, and let $T,t,p\in\mathbb{N}$. The input source list at time $t$ and position
$p$ consists, in order, of constants indicating whether $t=0$, $p=1$, $p=0$,
$p=|\ell|+1$, and $1\le p\le |\ell|$; the source $\ell_{p-1}$ when $1\le p\le |\ell|$,
and the false constant otherwise; the three input-history sources at positions
$p-1,p,p+1$ with their respective time and boundary flags; the two instruction-history
sources at position $p$ with indices $1$ and $2$; followed by the previous-snapshot
sources at time $t$.
\end{definition}

\begin{definition}[Snapshot source list]
\lean{Complexity.snapSrcs}
\label{Complexity.snapSrcs}
\leanok
\uses{BoolCircuit.BitSrc, Complexity.cellIdx, Complexity.inpIdx, Complexity.prevSnapSrcs, Complexity.snapIdx, Complexity.snapWidth}
Let $M$ be a finite Turing machine with Boolean symbols, let $\ell$ be a finite list of
bit sources, and let $T,t$ be natural numbers. The source list at time $t$ is, in order,
the constant indicating whether $t=0$; the constant false if $t=0$, and otherwise bit
$\mathrm{snapWidth}(M)$ of the snapshot block at time $t-1$; bits $1$ and $2$ of the
input block indexed by $M,\ell,T,t$ and $|\ell|+1$; for each of the $k$ machine states,
in order, bits $3$ and $4$ of the cell block indexed by $M,\ell,T,t$, that state, and
$2T$; and the previous-snapshot sources at time $t$.
\end{definition}

\begin{definition}[Padding a list of bit sources]
\lean{Complexity.padSrcs}
\label{Complexity.padSrcs}
\leanok
\uses{BoolCircuit.BitSrc, Complexity.cfgArity}
Let $M$ be a finite Turing machine with Boolean symbols, and let $l$ be a finite list of
bit sources. Define its padded list by appending to $l$ exactly $\text{cfgArity}(M)-|l|$
copies of the constant false source.
\end{definition}

\begin{definition}[Configuration instruction]
\lean{Complexity.cfgInstr}
\label{Complexity.cfgInstr}
\leanok
\uses{BoolCircuit.GInstr, Complexity.CfgKind, Complexity.cellSrcs, Complexity.cfgPos, Complexity.inpSrcs, Complexity.padSrcs, Complexity.snapSrcs}
Let $M$ be a finite Turing machine with Boolean symbols, let $\ell$ be a finite list of
bit sources, let $T\in\mathbb{N}$, and let $i\in\mathbb{N}$. The instruction at position
$i$ is determined by the configuration position assigned to $i$: for a cell position it
selects the machine-state kind associated with that state and uses the padded cell
sources; for an input position it selects the input kind and uses the padded input
sources; and for a snapshot position it selects the snapshot kind and uses the padded
snapshot sources.
\end{definition}

\begin{definition}[Configuration program]
\lean{Complexity.cfgProg}
\label{Complexity.cfgProg}
\leanok
\uses{BoolCircuit.GInstr, Complexity.CfgKind, Complexity.cfgInstr, Complexity.cfgStride}
Let $M$ be a finite Turing machine with Boolean symbols, let $\ell$ be a finite list of
bit sources, and let $T$ be a natural number. The configuration program is the list
obtained by applying the configuration-instruction construction for $M$, $\ell$, and $T$
at each successive index from $0$ up to, but not including, $(T+1)$ times the
configuration stride for $M$, the length of $\ell$, and $T$.
\end{definition}

\begin{theorem}[Length of the configuration program]
\lean{Complexity.length_cfgProg}
\label{Complexity.length_cfgProg}
\leanok
\uses{Complexity.cfgProg, Complexity.cfgStride}
\difficulty{1}
Let $M$ be a finite-state multi-tape Turing machine over the binary alphabet, with $k$
work tapes, let $\ell$ be a finite list of bit sources, and let $T$ be a natural number.
Set $S=k(2T+1)+|\ell|+3$; the configuration program formed with stride $S$ and time
horizon $T$ has length $(T+1)S$.
\end{theorem}

\begin{theorem}[Configuration program entries]
\lean{Complexity.cfgProg_getElem}
\label{Complexity.cfgProg_getElem}
\leanok
\uses{Complexity.cfgInstr, Complexity.cfgProg}
\difficulty{1}
Let $M$ be a finite-state Turing machine, let $\ell$ be a list of bit sources, and let
$T$ be a natural number. For every index $i$ within the configuration program, the entry
at index $i$ is the instruction constructed for index $i$.
\end{theorem}

\begin{definition}[Configuration expansion]
\lean{Complexity.cfgExp}
\label{Complexity.cfgExp}
\leanok
\uses{Complexity.cellExp, Complexity.cfgPos, Complexity.inpExp, Complexity.snapExp}
Let $M$ be a finite Turing machine with Boolean symbols, let $\ell$ be a finite list of
bit sources, let $T\in\mathbb{N}$, let $w$ be a Boolean list, and let $i\in\mathbb{N}$.
The configuration expansion at position $i$ is the cell expansion, input expansion, or
snapshot expansion specified by the configuration position assigned to $i$.
\end{definition}

\begin{theorem}[Configuration expansion at a work-cell index]
\lean{Complexity.cfgExp_cellIdx}
\label{Complexity.cfgExp_cellIdx}
\leanok
\uses{Complexity.cellExp, Complexity.cellIdx, Complexity.cfgExp, Complexity.cfgPos_cellIdx}
\difficulty{1}
Let $M$ be a finite-state binary Turing machine, let $\ell$ be a finite list of bit
sources, let $T,t\in\mathbb{N}$, let $\tau$ be one of its work tapes, and let $w$ be a
finite binary word. For every $r\in\mathbb{N}$ with $r<2T+1$, the configuration
expansion at the index assigned to time $t$, tape $\tau$, and cell coordinate $r$ equals
the work-cell expansion for $w,T,t,\tau,r$.
\end{theorem}

\begin{theorem}[Configuration expansion at an input index]
\lean{Complexity.cfgExp_inpIdx}
\label{Complexity.cfgExp_inpIdx}
\leanok
\uses{Complexity.cfgExp, Complexity.cfgPos_inpIdx, Complexity.inpExp, Complexity.inpIdx}
\difficulty{1}
Let $M$ be a finite-state binary Turing machine, let $\ell$ be a finite list of bit
sources, let $T,t\in\mathbb{N}$, and let $w$ be a finite binary word. For every
$p\in\mathbb{N}$ with $p<|\ell|+2$, the configuration expansion at the index assigned to
time $t$ and input position $p$ equals the input-position expansion for $M,w,t,p$.
\end{theorem}

\begin{theorem}[Configuration encoding at snapshot positions]
\lean{Complexity.cfgExp_snapIdx}
\label{Complexity.cfgExp_snapIdx}
\leanok
\uses{Complexity.cfgExp, Complexity.cfgPos_snapIdx, Complexity.snapExp, Complexity.snapIdx}
\difficulty{1}
Let $M$ be a finite-state binary Turing machine with $k$ work tapes, let $\ell$ be a
list of bit sources, and let $T,t\in\mathbb N$ and $w$ be a binary input. At the tableau
index for the snapshot at time $t$, the configuration encoding is exactly the snapshot
encoding at time $t$.
\end{theorem}

\begin{theorem}[Initial work-chain values]
\lean{Complexity.workChain_zero}
\label{Complexity.workChain_zero}
\leanok
\uses{Complexity.CfgTableau.cfgAt, Complexity.workChain, Turing.Cfg.init, Turing.Cfg.workTapeSymbols, Turing.MultiTapeTM.runFrom_zero}
\difficulty{3}
Let $M$ be a finite-state Turing machine with Boolean work-tape symbols, let $w$ be a
finite Boolean input, let $\tau$ be a work-tape index, and let $z\in\mathbb{Z}$. At time
zero, the work-chain value at position $z$ is none.
\end{theorem}

\begin{theorem}[Work chain vanishes to the left]
\lean{Complexity.workChain_left}
\label{Complexity.workChain_left}
\leanok
\uses{Complexity.abs_workTapePos_le, Complexity.workChain}
\difficulty{3}
Let $M$ be a finite-state Turing machine with Boolean work-tape symbols, let $w$ be a
finite Boolean input, and let $t,T\in\mathbb{N}$ with $t\le T$. For every work-tape
index $\tau$, the work-chain value at position $-T-1$ after $t$ steps is none.
\end{theorem}

\begin{theorem}[Work chain at the right boundary]
\lean{Complexity.workChain_right}
\label{Complexity.workChain_right}
\leanok
\uses{Complexity.CfgTableau.cfgAt, Complexity.abs_workTapePos_le, Complexity.workChain, Turing.Cfg.workTapeSymbols}
\difficulty{3}
Let $M$ be a finite-state Turing machine with Boolean work-tape symbols, let $w$ be a
finite Boolean input, and let $t,T\in\mathbb{N}$ with $t\le T$. For every work-tape
index $\tau$, the work-chain value at position $T$ after $t$ steps equals the symbol
scanned by tape $\tau$ at that time.
\end{theorem}

\begin{theorem}[Right boundary of the input chain]
\lean{Complexity.inpChain_right}
\label{Complexity.inpChain_right}
\leanok
\uses{Complexity.CfgTableau.cfgAt, Complexity.inpChain, Turing.Cfg.inputSymbol}
\difficulty{3}
Let $M$ be a finite-state multi-tape Turing machine over the binary alphabet, let $\ell$
be a list of circuit bit sources, let $T\in\mathbb N$, and let $w$ be a binary input.
For every $t\in\mathbb N$, the input-chain value at position $|w|+1$ after $t$ steps
equals the symbol read by the input head in the configuration reached after $t$ steps.
\end{theorem}

\begin{theorem}[Evaluation of an earlier block]
\lean{Complexity.eval_block}
\label{Complexity.eval_block}
\leanok
\uses{BoolCircuit.BitSrc.eval, Complexity.cfgExp}
\difficulty{3}
Let $M$ be a finite Turing machine over the Boolean alphabet, let $\ell$ be a list of
bit sources, and let $T$ be a time parameter. Let $E_w(r)$ denote the Boolean block
prescribed by the configuration encoding of input $w$ at index $r$. If a list of blocks
agrees with $E_w(r)$ at every index $r<i$, then, for any $a<i$ and any bit position $j$,
evaluating the block source that refers to block $a$, bit $j$ gives the $j$th bit of
$E_w(a)$, with false used when that bit position is out of range.
\end{theorem}

\begin{theorem}[Padding preserves the mapped prefix]
\lean{Complexity.padSrcs_map_take}
\label{Complexity.padSrcs_map_take}
\leanok
\uses{BoolCircuit.BitSrc, Complexity.padSrcs}
\difficulty{3}
Let $M$ be a finite-state multi-tape Turing machine over the binary alphabet, let $f$
map bit sources to Boolean values, and let $l$ be a finite list of bit sources. If $l$
is padded by appending false constants until its length is at least the arity associated
to $M$, then mapping $f$ over that padded list and taking its first $|l|$ entries gives
exactly the list obtained by mapping $f$ over $l$.
\end{theorem}

\begin{theorem}[Dropping and taking padded sources]
\lean{Complexity.drop_take_padSrcs}
\label{Complexity.drop_take_padSrcs}
\leanok
\uses{BoolCircuit.BitSrc, Complexity.padSrcs}
\difficulty{3}
Let $M$ be a finite Turing machine over the binary alphabet, let $f$ map bit sources to
Boolean values, and let $\mathrm{pre}$ and $\mathrm{rest}$ be lists of bit sources.
Mapping $f$ over the padded concatenation $\mathrm{pre}++\mathrm{rest}$, then dropping
$|\mathrm{pre}|$ entries and taking $|\mathrm{rest}|$ entries, gives exactly the list
obtained by mapping $f$ over $\mathrm{rest}$.
\end{theorem}

\begin{theorem}[Length of Previous Snapshot Sources]
\lean{Complexity.length_prevSnapSrcs}
\label{Complexity.length_prevSnapSrcs}
\leanok
\uses{Complexity.prevSnapSrcs, Complexity.snapWidth}
\difficulty{1}
Let $M$ be a finite-state Turing machine over the binary alphabet, let $\ell$ be a list
of bit sources, and let $T$ and $t$ be natural numbers. The list of bit sources defined
for the previous snapshot at time $t$ has length equal to the number of possible
snapshots of $M$.
\end{theorem}

\begin{theorem}[Evaluation of previous-snapshot sources]
\lean{Complexity.prevSnapSrcs_eval}
\label{Complexity.prevSnapSrcs_eval}
\leanok
\uses{BoolCircuit.BitSrc.eval, Complexity.cfgExp, Complexity.cfgExp_snapIdx, Complexity.eval_block, Complexity.prevSnapSrcs, Complexity.snapEncode, Complexity.snapExp, Complexity.snapIdx, Complexity.snapshotAt}
\difficulty{3}
Let $M$ be a finite binary Turing machine, let $\ell$ be a list of bit sources, and let
$T$ be a natural number. For binary input $x$, a list of Boolean blocks, a word $w$, and
natural numbers $i,t$, suppose every block at index $i'<i$ equals the
configuration-expansion block at index $i'$, and the snapshot index for time $t$ is less
than $i$. Then evaluating, on $x$ and the given blocks, the list of sources for the
snapshot immediately preceding time $t+1$ yields the encoding of the snapshot of $M$ on
$w$ at time $t$.
\end{theorem}

\begin{theorem}[Correctness of a Cell Block]
\lean{Complexity.CfgTableau.cell_block}
\label{Complexity.CfgTableau.cell_block}
\leanok
\uses{BoolCircuit.BitSrc.eval, Complexity.CfgTableau.optBits, Complexity.abs_workTapePos_le, Complexity.cellExp, Complexity.cellIdx, Complexity.cellIdx_lt_succ, Complexity.cellSrcs, Complexity.cfgExp, Complexity.cfgExp_cellIdx, Complexity.cfgF, Complexity.cfgF_cell_succ, Complexity.cfgF_cell_zero, Complexity.chSrc, Complexity.drop_take_padSrcs, Complexity.eval_block, Complexity.le_cellIdx, Complexity.length_prevSnapSrcs, Complexity.nbSrc, Complexity.padSrcs, Complexity.prevSnapSrcs, Complexity.prevSnapSrcs_eval, Complexity.snapDecode_snapEncode, Complexity.snapIdx_lt_succ, Complexity.workChain, Complexity.workChain_left, Complexity.workChain_zero}
\difficulty{6}
Let $M$ be a finite-state binary Turing machine with $k$ work tapes, let $\ell$ be a
list of bit sources, and let $T,N$ be the parameters determining the tableau layout. For
a binary input $w$, a block sequence, time $t\le T$, tape $\tau$, and cell index
$r<2T+1$, suppose that every block whose index is strictly less than the index assigned
to cell $(t,\tau,r)$ equals the expected tableau block at that index. Then evaluating
the padded local source list for this cell against the input and blocks, and applying
the cell update rule, yields exactly the expected encoding of cell $(t,\tau,r)$.
\end{theorem}
